\documentclass[10pt,a4paper]{amsart}
\usepackage{amsmath,amssymb,amsthm,graphicx}
\usepackage{verbatim}
\newenvironment{reference}{\comment}{\endcomment}
%\newenvironment{reference}{}{}
\newenvironment{slogan}{\comment}{\endcomment}
\newenvironment{history}{\comment}{\endcomment}
% For commutative diagrams we use Xy-pic
\usepackage[all]{xy}

% We use 2cell for 2-commutative diagrams.
\xyoption{2cell}
% Theorem environments.
%
\theoremstyle{plain}
\newtheorem{theorem}[subsection]{Theorem}
\newtheorem{proposition}[subsection]{Proposition}
\newtheorem{lemma}[subsection]{Lemma}

\theoremstyle{definition}
\newtheorem{definition}[subsection]{Definition}
\newtheorem{example}[subsection]{Example}
\newtheorem{exercise}[subsection]{Exercise}
\newtheorem{situation}[subsection]{Situation}

\theoremstyle{remark}
\newtheorem{remark}[subsection]{Remark}
\newtheorem{remarks}[subsection]{Remarks}

\numberwithin{equation}{subsection}

% Macros
%
\def\lim{\mathop{\mathrm{lim}}\nolimits}
\def\colim{\mathop{\mathrm{colim}}\nolimits}
\def\Spec{\mathop{\mathrm{Spec}}}
\def\Hom{\mathop{\mathrm{Hom}}\nolimits}
\def\Ext{\mathop{\mathrm{Ext}}\nolimits}
\def\SheafHom{\mathop{\mathcal{H}\!\mathit{om}}\nolimits}
\def\SheafExt{\mathop{\mathcal{E}\!\mathit{xt}}\nolimits}
\def\Sch{\mathit{Sch}}
\def\Mor{\mathop{\mathrm{Mor}}\nolimits}
\def\Ob{\mathop{\mathrm{Ob}}\nolimits}
\def\Sh{\mathop{\mathit{Sh}}\nolimits}
\def\NL{\mathop{N\!L}\nolimits}
\def\CH{\mathop{\mathrm{CH}}\nolimits}
\def\proetale{{pro\text{-}\acute{e}tale}}
\def\etale{{\acute{e}tale}}
\def\QCoh{\mathit{QCoh}}
\def\Ker{\mathop{\mathrm{Ker}}}
\def\Im{\mathop{\mathrm{Im}}}
\def\Coker{\mathop{\mathrm{Coker}}}
\def\Coim{\mathop{\mathrm{Coim}}}

% Boxtimes
%
\DeclareMathSymbol{\boxtimes}{\mathbin}{AMSa}{"02}

%
% Macros for moduli stacks/spaces
%
\def\QCohstack{\mathcal{QC}\!\mathit{oh}}
\def\Cohstack{\mathcal{C}\!\mathit{oh}}
\def\Spacesstack{\mathcal{S}\!\mathit{paces}}
\def\Quotfunctor{\mathrm{Quot}}
\def\Hilbfunctor{\mathrm{Hilb}}
\def\Curvesstack{\mathcal{C}\!\mathit{urves}}
\def\Polarizedstack{\mathcal{P}\!\mathit{olarized}}
\def\Complexesstack{\mathcal{C}\!\mathit{omplexes}}
% \Pic is the operator that assigns to X its picard group, usage \Pic(X)
% \Picardstack_{X/B} denotes the Picard stack of X over B
% \Picardfunctor_{X/B} denotes the Picard functor of X over B
\def\Pic{\mathop{\mathrm{Pic}}\nolimits}
\def\Picardstack{\mathcal{P}\!\mathit{ic}}
\def\Picardfunctor{\mathrm{Pic}}
\def\Deformationcategory{\mathcal{D}\!\mathit{ef}}
\usepackage{cleveref}
\crefname{originallemma}{original lemma}{original lemmas}
\Crefname{originallemma}{Original lemma}{Original lemmas}
\makeatletter
\newlabel{lemma-coherent-filter}{{07UR}{source}{}{sourcepointer-1}{}}
\newlabel{lemma-coherent-filter@cref}{{[originallemma][0][]07UR}{[1][0][]source}{}{}{}}
\makeatother
\crefname{originallemma}{original lemma}{original lemmas}
\Crefname{originallemma}{Original lemma}{Original lemmas}
\makeatletter
\newlabel{lemma-prepare-filter-irreducible}{{07UQ}{source}{}{sourcepointer-2}{}}
\newlabel{lemma-prepare-filter-irreducible@cref}{{[originallemma][0][]07UQ}{[1][0][]source}{}{}{}}
\makeatother
\crefname{originallemma}{original lemma}{original lemmas}
\Crefname{originallemma}{Original lemma}{Original lemmas}
\makeatletter
\newlabel{lemma-property-initial}{{07US}{source}{}{sourcepointer-3}{}}
\newlabel{lemma-property-initial@cref}{{[originallemma][0][]07US}{[1][0][]source}{}{}{}}
\makeatother
\title{A variant of the criterion for properties of coherent sheaves, testing on sheaves with prescribed scheme theoretic support}
\author{The Stacks Project Authors}
\date{stacks-project:spaces-cohomology.tex:local-snapshot-20261012; GNU FDL 1.2 or later}
\begin{document}
\maketitle
\noindent\textit{Editorial note.} All delivered mathematical text is the author's own, copied from the pinned Stacks Project chapter file spaces-cohomology.tex; no statement or proof step was rewritten, reordered or removed.\par\medskip
\noindent\textit{Editorial note.} Preamble provenance: the pinned chapter file opens with \textbackslash{}input\{preamble\} and carries no preamble text of its own, so the project's single Stacks-wide preamble staged at sources/stacks-spaces-cohomology/preamble.tex is excerpted verbatim as three spans, lines 8-12 (verbatim package and the project's own reference/slogan/history display environments), lines 14-18 (Xy-pic setup) and lines 154-217 (theorem declarations and mathematical macros). The wrapper class and the amsmath/amssymb/amsthm/graphicx base packages are supplied by the assembler. No macro was invented or redefined.\par\medskip
\noindent\textit{Editorial note.} Same-chapter pointers are located at their \textbackslash{}label lines inside the pinned chapter file; cross-chapter pointers are located in the project's own tag list staged in the snapshot (sources/stacks-chow/tags.txt), at the verbatim lines carrying their canonical tags.\par\medskip
\noindent\textit{Allowed standard prerequisites (editorial).}\begin{itemize}
\item The initial criterion reducing a property of coherent sheaves to two conditions on the ideals of reduced irreducible subspaces. Same chapter spaces-cohomology.tex, label lemma-property-initial (invoked at proof line 2797). The proof only has to verify conditions (1) and (2) of that criterion; condition (1) is immediate and condition (2) is what the minimal counterexample argument establishes.
\item Preparation of a coherent module on a reduced irreducible subspace by a pushed forward ideal power. Same chapter spaces-cohomology.tex, label lemma-prepare-filter-irreducible (invoked at proof line 2817). It provides the injective map \$\textbackslash{}varphi : i\_*\textbackslash{}mathcal\{I\}\textasciicircum{}\{\textbackslash{}oplus r\} \textbackslash{}to \textbackslash{}mathcal\{G\}\$ used to compare the two sheaves.
\item The filtration of a coherent sheaf on a Noetherian space with subquotients pushed forward from ideals on irreducible supports. Same chapter spaces-cohomology.tex, label lemma-coherent-filter (invoked at proof line 2822). The filtration of \$\textbackslash{}Coker(\textbackslash{}varphi)\$ lets minimality of \$Z\$ together with condition (1) give property \$\textbackslash{}mathcal\{P\}\$ for that cokernel.
\end{itemize}\medskip
\ifdefined\hypertarget\hypertarget{sourcepointer-1}{}\fi
\noindent\textit{Source reference.} Original lemma 07UR: sources/stacks-spaces-cohomology/pinned-original-source.tex, source lines 2521--2521. Navigation pointer for \textbackslash{}ref\{lemma-coherent-filter\} used inside the counted span (line 2822); the argument that consumes it is retained in the span.\par\medskip
\ifdefined\hypertarget\hypertarget{sourcepointer-2}{}\fi
\noindent\textit{Source reference.} Original lemma 07UQ: sources/stacks-spaces-cohomology/pinned-original-source.tex, source lines 2478--2478. Navigation pointer for \textbackslash{}ref\{lemma-prepare-filter-irreducible\} used inside the counted span (line 2817); the argument that consumes it is retained in the span.\par\medskip
\ifdefined\hypertarget\hypertarget{sourcepointer-3}{}\fi
\noindent\textit{Source reference.} Original lemma 07US: sources/stacks-spaces-cohomology/pinned-original-source.tex, source lines 2616--2616. Navigation pointer for \textbackslash{}ref\{lemma-property-initial\} used inside the counted span (line 2797); the argument that consumes it is retained in the span.\par\medskip
\begin{lemma}
\label{lemma-property-higher-rank-cohomological-variant}
Let $S$ be a scheme. Let $X$ be a Noetherian algebraic space over $S$.
Let $\mathcal{P}$ be a property of coherent sheaves on $X$. Assume
\begin{enumerate}
\item For any short exact sequence of coherent sheaves on $X$
if two out of three have property $\mathcal{P}$ so does the third.
\item If $\mathcal{P}$ holds  for $\mathcal{F}^{\oplus r}$ for
some $r \geq 1$, then it holds for $\mathcal{F}$.
\item For every reduced closed subspace $i : Z \to X$ with
$|Z|$ irreducible there exists a coherent sheaf $\mathcal{G}$ on $X$
whose scheme theoretic support is $Z$ such that $\mathcal{P}$ holds for
$\mathcal{G}$.
\end{enumerate}
Then property $\mathcal{P}$ holds for every coherent sheaf on $X$.
\end{lemma}

\begin{proof}
We will show that conditions (1) and (2) of
Lemma \ref{lemma-property-initial} hold. This is clear for condition (1).
To show that (2) holds, let
$$
\mathcal{T} =
\left\{
\begin{matrix}
i : Z \to X \text{ reduced closed subspace with }|Z|\text{ irreducible such}\\
\text{ that }i_*\mathcal{I}\text{ does not have }\mathcal{P}
\text{ for some quasi-coherent }\mathcal{I} \subset \mathcal{O}_Z
\end{matrix}
\right\}
$$
If $\mathcal{T}$ is nonempty, then since $X$ is Noetherian, we can
find an $i : Z \to X$ which is minimal in $\mathcal{T}$. We will show
that this leads to a contradiction.

\medskip\noindent
Let $\mathcal{G}$ be the sheaf whose scheme theoretic support is $Z$ whose
existence is assumed in assumption (3). Let
$\varphi : i_*\mathcal{I}^{\oplus r} \to \mathcal{G}$ be as in
Lemma \ref{lemma-prepare-filter-irreducible}. Let
$$
0 = \mathcal{F}_0 \subset \mathcal{F}_1 \subset
\ldots \subset \mathcal{F}_m = \Coker(\varphi)
$$
be a filtration as in Lemma \ref{lemma-coherent-filter}. By minimality
of $Z$ and assumption (1) we see that $\Coker(\varphi)$ has
property $\mathcal{P}$. As $\varphi$ is injective we conclude using
assumption (1) once more that $i_*\mathcal{I}^{\oplus r}$ has property
$\mathcal{P}$. Using assumption (2) we conclude that $i_*\mathcal{I}$
has property $\mathcal{P}$.

\medskip\noindent
Finally, if $\mathcal{J} \subset \mathcal{O}_Z$ is a second quasi-coherent
sheaf of ideals, set $\mathcal{K} = \mathcal{I} \cap \mathcal{J}$
and consider the short exact sequences
$$
0 \to \mathcal{K} \to \mathcal{I} \to \mathcal{I}/\mathcal{K} \to 0
\quad
\text{and}
\quad
0 \to \mathcal{K} \to \mathcal{J} \to \mathcal{J}/\mathcal{K} \to 0
$$
Arguing as above, using the minimality of $Z$, we see that
$i_*\mathcal{I}/\mathcal{K}$ and $i_*\mathcal{J}/\mathcal{K}$
satisfy $\mathcal{P}$. Hence by assumption (1) we conclude that
$i_*\mathcal{K}$ and then $i_*\mathcal{J}$ satisfy $\mathcal{P}$.
In other words, $Z$ is not an element of $\mathcal{T}$ which is
the desired contradiction.
\end{proof}

\end{document}
